\lecture{10}
\title{Stochastic Gradient Descent and Estimation Error Bounding}

\begin{document}

\subsection{Stochastic Gradient Descent}

One issue with gradient descent is that
it's slow to compute when the function being minimized
is the average of $n$ loss functions across $n$ data points,
like so, because computing the gradient $\nabla f(w_t)$
becomes an $O(n)$ task.
\[ f(w) := \frac{1}{n} \sum_{i = 1}^n f_i(w). \]
Stochastic gradient descent sacrifices accuracy for speed like so:

\textbf{Definition.}
To compute $w_{t + 1}$ from $w_t$ in \textbf{stochastic gradient descent},
sample $i \sim \mathrm{Unif}\{1, \dots, n\}$
and compute $g_{t + 1} = \nabla f_i(w_t)$.
Then set $w_{t + 1} \leftarrow w_t - \eta_{t + 1} g_{t + 1}$;
the point is that $g_{t + 1}$ is an
unbiased estimate for $\nabla f(w_t)$,
since $\mathbb{E}[g_{t + 1}] = \nabla f(w_t)$.

\begin{remark}
    If you want to decrease the variance of the approximation
    $g_{t + 1} \approx \nabla f(w_t)$, you can also
    set $g_{t + 1}$ to be the average of $\nabla f_i(w_t)$
    across $b$ randomly-sampled indices $i$.
\end{remark}

Importantly, the loss $f(w_t)$ may increase
on some iterations of SGD,
even for arbitrarily small step sizes $\eta_t$.

We conclude our discussion on SGD
with the statement of a theorem, without proof.

\begin{theorem}{SGD Convergence}
    Suppose $f$ is convex with minimizer $w^*$ and satisfies
    $\|\nabla f(w)\| \leq B$ for all $w$.
    Suppose:
    \begin{itemize}
        \item We have $\|w^* - w_0\| \leq D$,
        where $w_0$ is the initial value for SGD.
        \item We have $\mathbb{E}[g_{t + 1} \mid w_t] = \nabla f(w_t)$
        for all $t$.
        \item We have $\mathbb{E}\left[\|g_t\|^2\right] \leq B^2$
        for all $t$.
    \end{itemize}
    Then SGD with step size $\eta_t := \frac{D}{B\sqrt{t}}$ yields:
    \[ 
    \mathbb{E}\left[
        f(\overline{w}_t) - f(w^*)
    \right] \leq DB \cdot \left[
        \frac{2 + \log t}{2\sqrt{t}}
    \right]
    ~~~ \text{ where } ~~~
    \overline{w}_t :=
    \frac{\sum_{s = 1}^t \eta_s w_{s - 1}}{\sum_{s = 1}^t \eta_s}.
    \]
\end{theorem}

\begin{remark}
    The 6.790 notes claim the RHS should have a coefficient of $DB^2$
    rather than $DB$, but this is incorrect
    by dimensional analysis alone.
    The details here, for 6.790's sake,
    are not too important regardless.
\end{remark}

\subsection{Estimation Error Bounding}

Recently, we've been restricting ourselves
to specific classes of hypotheses (regression, parametric models).
Let's consider machine learning
at a much more general level;
what can we say \emph{regardless} of
what the hypothesis class $\mathcal{H}$ looks like?

\textbf{Definition.}
For any hypothesis $h \in \mathcal{H}$,
there are two types of errors (or risks).
\begin{itemize}
    \item The \textbf{expected (test) error} of $h$ is
    $\mathcal{R}(h) := \mathbb{E}\left[\mathcal{L}(h(X), Y)\right]$.
    \item The \textbf{empirical (training) error} of $h$ is
    $\hat{\mathcal{R}}(h) := \frac{1}{n} \sum_{i = 1}^n
    \mathcal{L}(h(x_i), y_i)$.
\end{itemize}

Can we upper-bound the \textbf{estimation error}
between the expected error and the empirical error?

\begin{theorem}{Estimation Error Bound}
    Assume $\mathcal{L}(x, y) \in [0, 1]$
    and pick some $\delta \in (0, 1)$.
    Then with probability $1 - \delta$,
    the following is true
    for \emph{all} hypotheses $h \in \mathcal{H}$:
    \[
    \left|\hat{\mathcal{R}}(h) - \mathcal{R}(h)\right|
    \leq \sqrt{\frac{\log|\mathcal{H}| + \log(2/\delta)}{2n}}.
    \]
\end{theorem}

Note that the assumption $\mathcal{L}(x, y) \in [0, 1]$
makes sense for classification problems, say.
More generally, one can achieve a similar bound
up to a constant factor for models with bounded loss functions.

\begin{remark}
    It should make sense that the RHS grows as $|\mathcal{H}|$ grows.
    The reason why is that larger $|\mathcal{H}|$ yield
    higher rates of overfitting,
    in which case the empirical error
    becomes a less accurate estimation of the expected error.

    One might question what happens in the case of linear regression,
    where $\mathcal{H}$ is an \emph{infinitely-sized}
    hypothesis class, with one hypothesis for each of the
    infinitely many parameters $\theta \in \mathbb{R}^d$.
    The trick, then, is to replace $\mathbb{R}$
    with, say, the set of $64$-bit floating point numbers.
    This way, $|\mathcal{H}| \approx 2^{64d}$ and
    $\log |\mathcal{H}| \approx 64d \log 2$,
    and the RHS doesn't blow up.
\end{remark}

We'll now provide a proof of this theorem.
Expect the proof to be ``simplistic''---after all,
it has to work no matter what $\mathcal{H}$ is!

\begin{proof}
    The proof will rely on a statistical fact
    known as \textbf{Hoeffding's Inequality},
    which we state below without proof.
    \begin{quote}
        \textbf{Hoeffding's Inequality.}
        If $U_1, \dots, U_n \in [0, 1]$ are
        sampled (i.i.d.) from the same distribution $\mathbb{P}$, then:
        \[
        \mathbb{P}\left(
            \left|\bar{U}_n - \mu\right| \geq t
        \right) \leq 2 \exp\left(-2nt^2\right).
        \]
        Note that $\bar{U}_n$ is the average of the $\{U_i\}_{i = 1}^n$,
        and $\mu$ is the expected value $\mu = \mathbb{E}[U_i]$
        (the same for all $i$).
    \end{quote}
    Directly applying Hoeffding's Inequality
    with $U_i := \mathcal{L}(h(x_i), y_i)$ yields:
    \[ \mathbb{P}\left(
    \left|\hat{\mathcal{R}}(h) - \mathcal{R}(h)\right| \geq t
    \right) \leq 2\exp\left(-2nt^2\right). \]
    Then the union bound across all $h \in \mathcal{H}$ yields:
    \[ \mathbb{P}\left(
    \left|\hat{\mathcal{R}}(h) - \mathcal{R}(h)\right| \geq t
    ~ \text{ for some } h \in \mathcal{H}
    \right) \leq 2|\mathcal{H}|\exp\left(-2nt^2\right). \]
    Plugging in $t = \sqrt{\frac{\log|\mathcal{H}| + \log(2/\delta)}{2n}}$
    now directly yields the promised bound;
    you can check this in your head. ~ $\blacksquare$
\end{proof}

We conclude with one more similar-looking theorem.

\textbf{Definition.}
We establish the notation
$\hat{h} := \argmin_{h \in \mathcal{H}} \hat{\mathcal{R}}(h)$
and $h^* := \argmin_{h \in \mathcal{H}} \mathcal{R}(h)$.

Clearly $h^*$ is preferred to $\hat{h}$,
but training data and ERM can only get us to $\hat{h}$.
How do $h^*$ and $\hat{h}$ compare?

\begin{theorem}{ERM Risk Bound}
    Assuming $\mathcal{L}(x, y) \in [0, 1]$,
    for all $\delta \in (0, 1)$
    we have with probability $1 - \delta$ that:
    \[
    \mathcal{R}(\hat{h}) - \mathcal{R}(h^*)
    \leq 2 \cdot 
    \sqrt{\frac{\log |\mathcal{H}| + \log\left(2 / \delta\right)}{2n}}.
    \]
\end{theorem}

\begin{proof}
    The point is the following:
    \[ \mathcal{R}(\hat{h}) - \mathcal{R}(h^*)
    = \underbrace{\left[
        \mathcal{R}(\hat{h}) - \hat{\mathcal{R}}(\hat{h})
    \right]}_{\text{bound \#1}}
    + \underbrace{\left[
        \hat{\mathcal{R}}(\hat{h}) - \hat{\mathcal{R}}(h^*)
    \right]}_{\leq \ 0}
    + \underbrace{\left[
        \hat{\mathcal{R}}(h^*) - \mathcal{R}(h^*)
    \right]}_{\text{bound \#2}}. \]
    The RHS has two bracketed expressions marked
    ``bound \#1'' and ``bound \#2''
    that can each be upper-bounded using the
    ``Estimation Error Bound'' theorem.
    Summing these two upper bounds
    yields the promised result. ~ $\blacksquare$
\end{proof}

\end{document}